\chapter{Analyse et spécifications}
\label{chap:analyse}
Avant toute décision de conception, il convient de fixer le cadre auquel
l'appareil devra répondre. Ce chapitre établit d'abord le cahier des charges,
c'est-à-dire l'ensemble des exigences que le produit doit satisfaire, puis le
traduit en spécifications techniques servant de référence vérifiable.
Conformément à la démarche descendante annoncée en introduction, ces deux
éléments précèdent l'analyse fonctionnelle, présentée ensuite, qui décompose le
système en fonctions afin d'organiser la conception matérielle et logicielle. \\

\section{Cahier des charges}
Le sujet ayant été choisi librement, le cahier des charges ne découle pas d'une
commande externe, il a été établi en début de projet et tient lieu de contrat
que l'appareil doit honorer. Son caractère auto-imposé ne le rend pas moins
contraignant, il délimite le périmètre du produit et sert de référence à
l'ensemble des choix qui suivent.

L'appareil doit être un lecteur audio portable, autonome et compact, dédié à la
seule lecture de musique. Cette vocation d'objet à fonction unique,
transportable dans une poche et utilisable sans dépendre d'un autre appareil,
oriente l'ensemble des exigences ci-dessous.  \\

\subsection{Exigences minimales}
Sur le plan audio, le produit doit lire des fichiers aux formats \texttt{.MP3}
et \texttt{.WAV} stockés sur une carte microSD, et restituer le son en stéréo
selon deux voies au choix de l'utilisateur : une sortie filaire sur jack
3,5\,mm et une sortie sans fil Bluetooth. Le volume d'écoute doit être réglable.

Sur le plan énergétique, l'appareil doit fonctionner sur batterie lithium-ion
avec une autonomie minimale de dix heures, et se recharger par un connecteur
USB-C.

Sur le plan de l'interface, un écran OLED graphique doit afficher au minimum le
titre en cours de lecture et l'état de la batterie, et l'utilisateur doit
pouvoir naviguer dans un système de menus à l'aide de commandes physiques.

Enfin, ce même connecteur USB-C doit permettre la programmation du
microcontrôleur, afin de n'exposer qu'une seule prise sur le boîtier.  \\

\subsection{Objectifs supplémentaires}
Au-delà de ce exigences minimales, les fonctions suivantes sont
extensions souhaitables si le calendrier le permet :
\begin{itemize}
    \item mise à jour du firmware depuis la carte microSD, Bluetooth ou Wi-Fi.
    \item synchronisation de l'heure via Bluetooth.
    \item interface graphique avancée avec des animations, icônes et affichant en permanence l'heure, l'état de la
          batterie et la connectivité active.
    \item navigation dans la bibliothèque musicale, avec tri par artiste, album
          et titre.
    \item menus de réglages (général, Bluetooth, égaliseur).
    \item écran de statistiques (tension, courant, température, luminosité).  \\
\end{itemize}

\section{Spécifications techniques}
Là où le cahier des charges énonce les fonctions attendues, les spécifications
techniques en fixent les exigences vérifiables. Imposées en début de projet,
elles constituent la référence contre laquelle les performances de l'appareil
seront contrôlées au chapitre consacré aux tests et à la validation. Le
tableau~\ref{tab:specs} les récapitule.

\begin{table}[h]
    \centering
    \caption{Spécifications techniques minimales}
    \label{tab:specs}
    \begin{tabular}{lll}
        \hline
        \textbf{Réf.} & \textbf{Spécification} & \textbf{Exigence}                 \\
        \hline
        SP1           & Restitution audio      & Stéréo, 24\,bits / 192\,kHz       \\
        SP2           & Formats décodés        & \texttt{.MP3} et \texttt{.WAV}    \\
        SP3           & Sortie filaire         & Jack 3.5\,mm                      \\
        SP4           & Sortie sans fil        & Bluetooth A2DP                    \\
        SP5           & Affichage              & OLED graphique                    \\
        SP6           & Stockage               & Carte microSD                     \\
        SP7           & Autonomie              & $\geq$ 10\,h sur batterie Li-ion  \\
        SP8           & Alimentation externe   & Charge et programmation via USB-C \\
        \hline
    \end{tabular}
\end{table}

\clearpage
\section{Analyse fonctionnelle matérielle}
L'analyse fonctionnelle traduit les exigences du cahier des charges en une
décomposition du système, indépendamment du détail des composants. Elle est
conduite à deux niveaux de granularité : le niveau~0 considère l'appareil comme
une boîte noire et recense ses échanges avec l'environnement extérieur, le
niveau~1 ouvre cette boîte et répartit les fonctions internes en blocs reliés
par leurs flux d'énergie et de données. Cette décomposition sert ensuite de
trame à la conception matérielle. \\

\subsection{Niveau 0 : diagramme de contexte}
Le niveau~0, présenté à la figure~\ref{fig:af-niv0}, situe l'appareil dans son
environnement et identifie ses interfaces avec l'extérieur. L'appareil échange
avec quatre éléments externes :
\begin{itemize}
    \item Le \textbf{port USB-C}, par lequel transitent l'énergie de recharge
          et la programmation du microcontrôleur.
    \item La \textbf{carte microSD}, support amovible contenant les fichiers
          musicaux à lire.
    \item Le \textbf{jack 3.5\,mm}, sortie audio filaire vers un casque ou des
          écouteurs.
    \item La liaison \textbf{Bluetooth}, sortie audio sans fil vers un appareil
          d'écoute distant.
\end{itemize}
À l'exception de la liaison radio, qui rayonne directement par l'antenne,
chacune de ces interfaces physiques traverse un étage de protection à la
frontière du système. Cet étage, qui protège les broches exposées contre les
décharges électrostatiques et les perturbations, constitue la première
fonction rencontrée par tout signal entrant ou sortant. \\

\fig[H, width=0.7\textwidth]{Analyse fonctionnelle matérielle - Niveau 0}{Analyse_Fonctionnelle-Niv1.drawio}
\label{fig:af-niv0}

\clearpage
\subsection{Niveau 1 : décomposition interne}
Le niveau~1, présenté à la figure~\ref{fig:af-niv1}, décompose le système en
blocs fonctionnels et explicite les flux qui les relient. On distingue les flux
d'énergie, issus de l'alimentation, et les flux de données,
échangés autour du processeur.

\paragraph{Énergie :} Le port USB-C fournit le 5\,V d'entrée, qui alimente le
bloc \emph{Alimentation et batterie} après passage par l'étage de protection.
Ce bloc assure la charge de l'accumulateur lithium-ion et la génération du
3.3\,V régulé qui alimente l'ensemble de l'électronique, à commencer par le
processeur.

\paragraph{Traitement et données :} Le bloc \emph{Processeur et communication
    sans-fil} occupe la position centrale : il orchestre l'ensemble des échanges.
Il lit les fichiers musicaux sur la carte microSD via SPI, pilote
l'affichage via un autre bus SPI, interroge les boutons via l'expanseur
d'entrées-sorties via I\textsuperscript{2}C, et émet le flux audio sans
fil par son antenne Bluetooth intégrée. La programmation s'effectue depuis le
port USB-C par l'intermédiaire d'un pont USB-vers-UART.

\paragraph{Chaîne audio filaire :} Pour la sortie filaire, le processeur
transmet l'audio numérique au bloc \emph{Audio filaire} via
I\textsuperscript{2}S pour les échantillons et via I\textsuperscript{2}C pour la
configuration. Ce bloc convertit le fichier numérique en signal en analogique, l'amplifie, puis le
restitue en stéréo sur le jack 3.5\,mm après passage par l'étage de protection.

\fig[H, width=1\textwidth]{Analyse fonctionnelle matérielle - Niveau 1}{Analyse_Fonctionnelle-Niv2.drawio}
\label{fig:af-niv1}

Cette décomposition fonctionnelle découpe les sous-systèmes alimentation,
traitement, interface utilisateur, stockage et audio, qui structureront le
chapitre de conception matérielle, où chaque bloc est traduit en un choix de
composants et un schéma.
